%!TEX TS-program = lualatex
%!TEX encoding = UTF-8
% Gabarit d'un tableau périodique ombré : un style à argument reçoit, de 0 à 100, la
% valeur normalisée que produit TableauTikz.ombrage. Les cases, symboles et
% informations viennent de ombrage_chi.tex, que produit ombrage_chi.py :
%     python ombrage_chi.py && lualatex Tableau_ombrage_chi.tex
\documentclass{standalone}
\usepackage{tikz}
\usepackage{siunitx}

\newcommand\dessinecase[2][]{\draw[#1] (#2) ++(-.48, -.48) rectangle ++(.96, .96);}

\def\larg{1cm}
\def\haut{1cm}
\def\xoriginlanth{3.5}
\def\yoriginlanth{-1}

\sisetup{output-decimal-marker = {,}}

\begin{document}
\fontsize{9}{9}\selectfont

\begin{tikzpicture}[x=\larg, y=\haut,
	% le style à argument : #1 vaut de 0 (chi = 0) à 100 (chi = 4), bornes du script
	chi/.style={fill=red!#1!blue!60!white, draw=black!50},
	sans chi/.style={fill=black!10, draw=black!50, densely dotted},
	titre colonne/.style={gray, scale=.7, node font={\itshape\footnotesize}, yshift=-.4\haut},
	titre ligne/.style={gray, scale=.7, node font={\itshape\footnotesize}, xshift=.4\larg},
	symbole/.style={node font={\bfseries\normalsize}, yshift=.05\haut},
	chi info/.style={node font={\tiny}, yshift=-.3\haut},
	nom/.style={scale=.7, node font={\tiny}, yshift=.5\haut, text centered},
	Z/.style={node font={\bfseries\tiny}, yshift=.2\haut, xshift=.3\larg}
]

% coordonnées <groupe>-<période> du tableau principal
\foreach \col in {1,...,18} {
	\foreach \ligne in {1,...,7} {
		\draw (\col, 7-\ligne) coordinate(\col-\ligne);
	}
}

% titres
\foreach \col in {1,18} \draw (\col, 7) node[titre colonne]{\col};
\foreach \col in {2, 13, 14, 15, 16, 17} \draw (\col, 6) node[titre colonne]{\col};
\foreach \col in {3, ..., 12} \draw (\col, 4) node[titre colonne]{\col};
\foreach \ligne in {1,...,7} \draw (0, {(7-\ligne)}) node[titre ligne]{\ligne};

% coordonnées <rang>-L et <rang>-A des lignes du bloc f
\foreach \col in {1,...,14} {
	\draw (\xoriginlanth-1, \yoriginlanth) ++ (\col, 0) node[titre colonne]{\col};
	\draw ({\xoriginlanth+\col-1}, {\yoriginlanth-1}) coordinate(\col-L);
	\draw ({\xoriginlanth+\col-1}, {\yoriginlanth-2}) coordinate(\col-A);
}

% éléments
\input{ombrage_chi}

% renvoi des lignes du bloc f, entre les groupes 2 et 3
\draw[very thick] (2.5, -.5) -- ++(0,2) (\xoriginlanth-.5, \yoriginlanth-.5) -- ++(0,-2);
\draw (2.5, -.5) |- (\xoriginlanth-.5, \yoriginlanth-1.5);

% légende : la même échelle que le style, de chi = 0 (bleu) à chi = 4 (rouge)
\shade[left color=red!0!blue!60!white, right color=red!100!blue!60!white, draw=black!50]
	(4, 5.2) rectangle ++(6, .4);
\foreach \x/\chi in {0/0, 1.5/1, 3/2, 4.5/3, 6/4} {
	\draw (4+\x, 5.2) -- ++(0, -.1) node[below, font=\footnotesize] {\num{\chi}};
}
\draw (7, 5.75) node[font=\footnotesize] {électronégativité de Pauling $\chi$ ; en pointillé : valeur inconnue};
\end{tikzpicture}

\end{document}
